\chapter{Strictly Richer Physics at Each Stage}
\label{sec:richer}
% ===================================================================

We now argue that the physics --- in the full sense developed in
Part~I --- is strictly richer at each stage of the fibration.
This is not merely a statement about computational power;
it is a statement about ontological resolution.

\section{Finer Bisimulation Quotient}

\begin{proposition}[Strict Refinement of Bisimulation]
\label{prop:strict-bisim}
For any $\GSLT \in \mathbf{GSLT}_{\mathrm{TC}}$ and any ordinal $\alpha$,
the bisimulation quotient at stage $\alpha+1$ is strictly finer than
at stage $\alpha$:
\[
  {\bisim_{\alpha+1}} \;\subsetneq\; {\bisim_\alpha}
\]
as equivalence relations on $\terms(\fiber{\alpha+1})$.
That is, there exist terms $P, Q$ such that $P \bisim_\alpha Q$ but
$P \not\bisim_{\alpha+1} Q$.
\end{proposition}

\begin{proof}[Proof sketch]
By a diagonalization argument.
Consider any pair of terms $P, Q$ whose bisimilarity in
$\fiber{\alpha}$ depends on whether some computation in
$\fiber{\alpha}$ halts.
At stage $\alpha+1$, the oracle can settle this halting question,
distinguishing $P$ from $Q$ in one interaction step.
Therefore $P \not\bisim_{\alpha+1} Q$ while $P \bisim_\alpha Q$.
\end{proof}

\begin{remark}
This means the \emph{atoms of observable behavior} are different
at each stage.
Terms that are experimentally indistinguishable at stage $\alpha$
--- that no agent at stage $\alpha$ can tell apart, no matter how
long it experiments --- become distinguishable at stage $\alpha+1$.
The world has more things in it at the higher stage.
\end{remark}

\section{Strictly More Expressive Hypothesis Language}

\begin{proposition}[Strict Expansion of HML]
The context-decorated HML at stage $\alpha+1$ strictly extends
the HML at stage $\alpha$:
\[
  \HML(\ctx_\alpha) \;\subsetneq\; \HML(\ctx_{\alpha+1}).
\]
In particular, there exist formulae $\phi \in \HML(\ctx_{\alpha+1})$
that are not expressible in $\HML(\ctx_\alpha)$.
\end{proposition}

\begin{proof}[Proof sketch]
The oracle interaction gives rise to new minimal contexts at stage
$\alpha+1$: contexts that place a term in interaction with the
oracle.
These generate new modal operators $\langle K_\oracle \rangle$
not present at stage $\alpha$.
Formulae using these operators are inexpressible at stage $\alpha$.
\end{proof}

\begin{remark}
An agent at stage $\alpha$ not only cannot \emph{test} certain
hypotheses; it cannot even \emph{form} them.
The concepts required to think the thought are absent from the
hypothesis language.
This is the sense in which a higher-stage agent inhabits a
categorically richer world: it has concepts unavailable at lower
stages, not merely faster procedures.
\end{remark}

\section{New Conserved Quantities}

\begin{proposition}[New Noether Currents at Each Stage]
\label{prop:new-noether}
The weight and cost maps at stage $\alpha+1$ can be sensitive to
distinctions invisible at stage $\alpha$.
Symmetries of the stage-$(\alpha+1)$ cost map that are not
symmetries of the stage-$\alpha$ cost map give rise to new
conserved quantities at stage $\alpha+1$ that have no analogue
at stage $\alpha$.
\end{proposition}

\begin{remark}
In physical terms: the higher-stage physics has strictly more
conservation laws.
What appears at stage $\alpha$ as a single undifferentiated
interaction resolves at stage $\alpha+1$ into a structured process
with internal symmetries and corresponding conserved currents.
The analogy with the hierarchy of effective field theories is
direct: new conservation laws (baryon number, lepton number,
color charge) appear as one descends to finer scales.
\end{remark}

\section{Richer Internal Causal Structure}

\begin{proposition}[Resolution of Elementary Interactions]
\label{prop:resolution}
A transition $P \rewrite_\alpha Q$ that is a single step in
$\fiber{\alpha}$ may resolve, in $\fiber{\alpha+1}$, into a
structured multi-step causal history
$P = R_0 \rewrite_{\alpha+1} R_1 \rewrite_{\alpha+1} \cdots
\rewrite_{\alpha+1} R_k = Q$
with nontrivial internal causal structure visible only from
stage $\alpha+1$.
\end{proposition}

\begin{remark}
What appeared to be an elementary interaction at stage $\alpha$ ---
an atom of causal history --- has internal structure at stage
$\alpha+1$.
This is the precise sense in which the world is richer at higher
stages: not just that there are more things, but that the things
that existed before turn out to be composite.
The analogy with the discovery that atoms have internal structure
(electrons, nuclei) and that nuclei have internal structure
(protons, neutrons) and that protons have internal structure (quarks)
is direct and intentional.
\end{remark}

% ===================================================================
